% ===================================================================
\chapter{Background}
\label{ch:background}

This chapter covers what the radar measures and how, why its data is large,
how lossless compression works, and the FPGA concepts the hardware chapters
use.

\section{How an FMCW radar measures range}
\label{sec:fmcw}

\subsection{The chirp}

Automotive radar is almost universally \emph{frequency-modulated continuous
wave} (FMCW). Instead of emitting a short pulse and timing the echo --- which
would demand impractically fast electronics for the distances involved --- an
FMCW radar transmits continuously, sweeping its frequency linearly upward over
a short interval. One such sweep is called a \emph{chirp}. The sensor used here
sweeps \SI{1.798}{\giga\hertz} of bandwidth around a \SI{77}{\giga\hertz}
centre frequency while its echo is sampled, and starts a new chirp every
\SI{78.5}{\micro\second}.

When the chirp reflects off an object at range $R$, the echo arrives delayed by
the round-trip time
\begin{equation}
\tau = \frac{2R}{c},
\end{equation}
where $c$ is the speed of light. Because the transmitter has kept sweeping in
the meantime, the echo's frequency at any instant is lower than what is
currently being transmitted, by an amount proportional to the delay. If the
chirp sweeps $B$ hertz in $T_c$ seconds, the sweep slope is $S = B/T_c$ and the
frequency difference is
\begin{equation}
\label{eq:beat}
f_b = S\,\tau = \frac{2 S R}{c}.
\end{equation}

\begin{figure}[t]
\centering
\begin{tikzpicture}[font=\small]
% axes
\draw[->,slate] (0,0) -- (10.8,0) node[right,font=\scriptsize]{time};
\draw[->,slate] (0,0) -- (0,4.3) node[above,font=\scriptsize]{frequency};
% legend (clear of the chirps)
\draw[accent,line width=1.2pt] (0.45,4.0) -- (1.15,4.0);
\node[font=\scriptsize,text=accent,anchor=west] at (1.25,4.0) {transmitted chirp};
\draw[warm,line width=1.2pt,densely dashed] (4.3,4.0) -- (5.0,4.0);
\node[font=\scriptsize,text=warm,anchor=west] at (5.1,4.0) {received echo (delayed by $\tau$)};
% transmitted chirps
\foreach \s in {0,3.2,6.4}{ \draw[accent,line width=1.2pt] (\s,0.35) -- (\s+2.6,2.9); }
% received (delayed)
\foreach \s in {0,3.2,6.4}{ \draw[warm,line width=1.2pt,densely dashed] (\s+0.55,0.35) -- (\s+3.15,2.9); }
% delay + beat annotation on first chirp
\draw[slate,dotted] (1.25,0.87) -- (2.05,0.87);
\draw[<->,slate,line width=0.5pt] (1.25,0.87) -- (2.05,0.87)
      node[midway,below,font=\scriptsize,text=slate,yshift=-1pt]{$\tau$};
\draw[<->,slate,line width=0.5pt] (2.05,0.87) -- (2.05,1.65)
      node[midway,right,font=\scriptsize,text=slate,xshift=1pt]{$f_b$};
% chirp duration brace
\draw[decorate,decoration={brace,amplitude=4pt,mirror},slate]
   (0,-0.28) -- (2.6,-0.28) node[midway,below,font=\scriptsize,text=slate,yshift=-3pt]{$T_c$};
\draw[decorate,decoration={brace,amplitude=4pt,mirror},slate]
   (0,-1.15) -- (9.0,-1.15) node[midway,below,font=\scriptsize,text=slate,yshift=-3pt]
   {one frame: many chirps --- the \emph{slow time} axis};
\end{tikzpicture}
\caption[FMCW chirp and beat frequency]{FMCW operation. Each chirp sweeps
linearly in frequency. The echo is a delayed copy, so at any instant the
transmitted and received frequencies differ by $f_b = S\tau$, proportional to
target range. A frame consists of many such chirps in succession; the axis
running across chirps is called \emph{slow time} and is where Doppler
information lives.}
\label{fig:chirp}
\end{figure}

A mixer multiplies the transmitted and received signals, producing a
low-frequency \emph{beat} or intermediate-frequency signal at exactly $f_b$.
That signal is low-pass filtered and digitised by an analog-to-digital
converter (ADC) at a modest rate --- \SI{8}{\mega\hertz} here, rather than the
tens of gigahertz the raw RF would need. This is the central trick of FMCW: it
converts a timing measurement that would need picosecond resolution into a
frequency measurement that an ordinary ADC can make.

\subsection{From beat frequency to a range profile}

Measuring $f_b$ is a spectral estimation problem, so the natural tool is the
Fourier transform. Taking an $N_s$-point FFT across the $N_s$ ADC samples of a
single chirp produces a \emph{range profile}: bin $n$ of the FFT corresponds,
through Equation~\eqref{eq:beat}, to range
\begin{equation}
R_n = n \cdot \frac{c}{2B},
\end{equation}
so the range resolution is
\begin{equation}
\Delta R = \frac{c}{2B} \approx \SI{8.3}{\centi\metre}
\end{equation}
for the \SI{1.798}{\giga\hertz} bandwidth used here. A peak in the range
profile is an object. This first transform is called the \textbf{Range FFT},
and its output is the data this thesis compresses.

\subsection{Slow time and Doppler}

A single chirp gives range but not velocity. To get velocity, the radar
transmits many chirps in rapid succession --- a \emph{frame} --- and looks at
how the \emph{phase} of a given range bin evolves from chirp to chirp.

The reasoning is worth spelling out, because the whole of DRHE rests on it. A
target at range $R$ produces a range-FFT value whose phase is
$\theta = -4\pi R/\lambda$, where $\lambda$ is the wavelength. Over the
chirp-to-chirp interval $T$, a target moving at radial velocity $v$ changes
range by $vT$, so its phase changes by
\begin{equation}
\label{eq:dopplerphase}
\Delta\theta = -\frac{4\pi v T}{\lambda}.
\end{equation}
At \SI{77}{\giga\hertz}, $\lambda \approx \SI{3.9}{\milli\metre}$, so even a
very slow target produces a measurable phase step --- which is exactly why
radar is good at velocity. Critically for compression, $\Delta\theta$ is
\emph{constant} for a target at constant velocity. The phase of a given range
bin therefore advances linearly across chirps, and its magnitude barely changes
at all. That is a highly predictable sequence, and predictable sequences
compress.

The axis running across chirps within a frame is called \textbf{slow time}, as
opposed to the \emph{fast time} axis of samples within a chirp. Taking a second
FFT along slow time resolves $\Delta\theta$ into velocity: this is the
\textbf{Doppler FFT}.

\section{Angle, MIMO and the virtual array}
\label{sec:mimo}

Angle is measured by comparing the phase of the same echo at separated receive
antennas: a wavefront arriving at angle $\phi$ reaches adjacent antennas,
spaced $d$ apart, with a phase difference
\begin{equation}
\label{eq:steering}
\Delta\psi = \frac{2\pi d \sin\phi}{\lambda},
\end{equation}
and an FFT across the antennas --- the \textbf{Angle FFT} --- turns these into
a bearing. Angular resolution grows with the size of the array.
Multiple-Input Multiple-Output (MIMO) radar enlarges the array cheaply: if the
signal of each of $N_t$ transmitters can be separated at each of $N_r$
receivers, every transmit--receive pair acts as an element of a \emph{virtual
array}. The sensor used here has $N_t = 12$ and $N_r = 16$, giving 192 virtual
elements from 28 physical antennas.

\subsection{Time-division multiplexing, and why it matters here}
\label{sec:tdm}

The separability requirement has to be met somehow. The simplest method, and
the one this sensor uses, is \emph{time-division multiplexing} (TDM): the
transmitters fire one at a time, in a fixed rotation, so whatever a receiver
hears during transmitter $t$'s slot must have come from transmitter $t$.

TDM is simple and robust, and it has one consequence that matters throughout
this thesis. Inside every repetition of the frame --- every ``chirp loop''
--- all $N_t$ transmitters fire in sequence before the loop repeats. When the
resulting data is laid out as a two-dimensional array of (ramp, range bin),
which is how a streaming compressor sees it, the ramp index is
\begin{equation}
\label{eq:rampindex}
m = q\,N_t + t,
\qquad q \in [0, N_c), \quad t \in [0, N_t),
\end{equation}
where $q$ indexes the chirp loop and $t$ the transmitter. The transmitter index
varies \emph{fastest}.

This is the correct physical time order. It is also, as Figure~\ref{fig:tdm}
makes visual, a trap. Ramp $m$ and ramp $m-1$ are different physical
transmitters $N_t - 1$ times out of $N_t$. They are separated in \emph{space}
by the array baseline, not in \emph{time} by a pulse interval, and the phase
relationship between them is governed by Equation~\eqref{eq:steering} --- the
target's bearing --- not by Equation~\eqref{eq:dopplerphase} --- the target's
velocity. The genuine slow-time neighbour of ramp $m$, the one whose phase
differs by a Doppler step, is ramp $m - N_t$. From Chapter~\ref{ch:methods} on,
this distance is written $\nTx$; on this sensor $\nTx = N_t = 12$.

\begin{figure}[t]
\centering
\begin{tikzpicture}[font=\small,x=0.50cm,y=0.50cm]
\foreach \i in {0,...,25}{
  \pgfmathtruncatemacro{\tx}{mod(\i,12)}
  \pgfmathsetmacro{\shade}{18+\tx*6}
  \node[draw=slate!55,line width=0.3pt,fill=accent!\shade,minimum width=0.46cm,
        minimum height=0.56cm,inner sep=0pt,font=\tiny] (c\i) at (\i,0) {\tx};
}
\node[anchor=east,font=\scriptsize,text=slate] at (-0.55,0) {TX index};
\node[anchor=west,font=\scriptsize,text=slate] at (26.0,0) {$\cdots$};
% highlight the three cells that matter
\node[draw=green2,line width=1.1pt,minimum width=0.46cm,minimum height=0.56cm,inner sep=0pt] at (2,0) {};
\node[draw=warm,line width=1.1pt,minimum width=0.46cm,minimum height=0.56cm,inner sep=0pt] at (13,0) {};
\node[draw=slate,line width=1.4pt,minimum width=0.46cm,minimum height=0.56cm,inner sep=0pt] at (14,0) {};
\node[font=\scriptsize,text=slate,anchor=south] at (14,0.40) {ramp $m$};
% ramp ruler
\foreach \i/\l in {0/0,12/12,24/24}{ \node[font=\tiny,text=slate] at (\i,-0.62) {\l}; }
\node[anchor=east,font=\scriptsize,text=slate] at (-0.55,-0.62) {ramp index};
% chirp loop braces
\draw[decorate,decoration={brace,amplitude=3pt},slate] (-0.3,1.55) -- (11.3,1.55)
   node[midway,above,font=\scriptsize,text=slate,yshift=-1pt]{chirp loop 0};
\draw[decorate,decoration={brace,amplitude=3pt},slate] (11.7,1.55) -- (23.3,1.55)
   node[midway,above,font=\scriptsize,text=slate,yshift=-1pt]{chirp loop 1};
% lag arrows, drawn under the strip
\draw[warm,line width=1.0pt,-{Latex[length=1.8mm]}] (14,-1.05) -- (13.1,-1.05);
\node[font=\scriptsize,text=warm,anchor=west] at (-0.3,-1.85)
   {\textbf{lag 1} lands on TX 1 --- a \emph{different antenna}; the phase step is the target's \emph{bearing}};
\draw[green2,line width=1.0pt,-{Latex[length=1.8mm]}] (14,-2.75) -- (2.1,-2.75);
\node[font=\scriptsize,text=green2,anchor=west] at (-0.3,-3.55)
   {\textbf{lag $n_{\mathrm{Tx}}\!=\!12$} lands on TX 2 --- the \emph{same antenna}; the phase step is the target's \emph{velocity}};
\end{tikzpicture}
\caption[TDM-MIMO ramp interleaving]{Why ``the previous ramp'' is the wrong
neighbour on a TDM-MIMO cascade. Each cell is one ramp, labelled with the
transmitter that produced it; the twelve transmitters rotate inside every chirp
loop. Predicting ramp $m$ from ramp $m-1$ (red) crosses a transmitter boundary
eleven times out of twelve, so the two ramps differ by the target's array
steering phase. Predicting from ramp $m - n_{\mathrm{Tx}}$ (green) stays on the
same physical antenna and therefore on the genuine Doppler axis.}
\label{fig:tdm}
\end{figure}


\section{The radar data cube, and why it is large}
\label{sec:cube}

Putting the three axes together gives the \emph{radar data cube}: fast time
(samples within a chirp) $\times$ slow time (ramps) $\times$ receive channel.
Three FFTs are applied along the three axes in turn, yielding range, velocity
and angle. Figure~\ref{fig:cube} shows the arrangement and where in the chain
this thesis intervenes.

\begin{figure}[t]
\centering
\begin{tikzpicture}[font=\small]
% cube
\begin{scope}[shift={(0,0)},scale=1.0]
  \draw[fill=lightbg,draw=slate] (0,0) -- (2.6,0) -- (2.6,2.0) -- (0,2.0) -- cycle;
  \draw[fill=accent!12,draw=slate] (2.6,0) -- (3.5,0.7) -- (3.5,2.7) -- (2.6,2.0) -- cycle;
  \draw[fill=accent!22,draw=slate] (0,2.0) -- (2.6,2.0) -- (3.5,2.7) -- (0.9,2.7) -- cycle;
  \node[font=\scriptsize,text=slate] at (1.3,-0.35) {fast time ($N_s$)};
  \node[font=\scriptsize,text=slate,rotate=90] at (-0.35,1.0) {slow time ($N_c N_t$)};
  \node[font=\scriptsize,text=slate,rotate=38] at (3.35,1.55) {RX ($N_r$)};
\end{scope}
% chain
\begin{scope}[shift={(5.1,0.25)},
  b/.style={draw=slate,fill=white,rounded corners=2pt,minimum width=19mm,minimum height=7.5mm,font=\scriptsize},
  h/.style={draw=accent,line width=1pt,fill=lightbg,rounded corners=2pt,minimum width=19mm,minimum height=7.5mm,font=\scriptsize}]
  \node[b] (a) at (0,2.1) {ADC samples};
  \node[b] (b) at (0,1.05) {window + DC removal};
  \node[b] (c) at (0,0.0) {\textbf{Range FFT}};
  \node[h] (d) at (0,-1.05) {FX16 quantise};
  \node[h] (e) at (0,-2.1) {\textbf{compress}};
  \node[b] (f) at (4.4,-2.1) {transmit / store};
  \node[h] (g) at (4.4,-1.05) {decompress};
  \node[b] (h2) at (4.4,0.0) {Doppler FFT};
  \node[b] (i) at (4.4,1.05) {Angle FFT};
  \node[b] (j) at (4.4,2.1) {detection (CFAR)};
  \foreach \x/\y in {a/b,b/c,c/d,d/e} {\draw[-{Latex[length=1.5mm]},slate] (\x) -- (\y);}
  \draw[-{Latex[length=1.5mm]},slate] (e) -- (f);
  \foreach \x/\y in {f/g,g/h2,h2/i,i/j} {\draw[-{Latex[length=1.5mm]},slate] (\x) -- (\y);}
  \node[font=\scriptsize,text=accent,align=left,anchor=west] at (-1.5,-3.0)
     {\textbf{scope of this thesis}};
  \draw[accent,dashed,rounded corners=3pt] (-1.25,-2.6) rectangle (5.65,-0.6);
\end{scope}
\end{tikzpicture}
\caption[Radar data cube and processing chain]{The radar data cube and the
processing chain. The compressor sits immediately after range-FFT and
fixed-point quantisation, before the data leaves the sensor. Everything
downstream --- Doppler FFT, angle FFT, detection --- sees the reconstructed
data, which for a lossless compressor is bit-identical to the original.}
\label{fig:cube}
\end{figure}

The volume follows directly. For the sensor used here, one frame is
\begin{equation}
N_s \times N_c \times N_r \times N_t
 = 256 \times 16 \times 16 \times 12 = \num{786432}
\end{equation}
complex samples. Each is two 16-bit integers, so a frame is
\SI{3145728}{\byte} --- exactly \SI{3}{\mebi\byte}. At 20 frames per second,
a common automotive rate, that is about half a gigabit per second from one
sensor, and a vehicle carries several.

\section{Lossless compression: how it actually works}
\label{sec:compression}

\subsection{The two categories}

Compression divides into \emph{lossless}, where the decompressed data is
bit-identical to the input, and \emph{lossy}, where some information is
permanently discarded in exchange for a smaller payload \cite{sayood2017}.
Lossy compression can reach far higher ratios; lossless compression is bounded
by the information content of the data itself. Section~\ref{sec:lossless}
explained why this application needs the lossless kind.

\subsection{The bound: entropy}

There is a hard limit on lossless compression, and knowing it is what makes the
results in Chapter~\ref{ch:results} interpretable. If a source emits symbols
$x$ with probabilities $p(x)$, its \emph{order-0 entropy} is
\begin{equation}
\label{eq:entropy}
H = -\sum_x p(x) \log_2 p(x) \quad\text{bits per symbol},
\end{equation}
and no code that treats symbols independently can average fewer bits than $H$
\cite{sayood2017}. This gives a concrete yardstick: measure $H$ for the data,
compare it against what the coder actually achieves, and the gap is the coder's
inefficiency. Section~\ref{sec:decomp} uses exactly this to separate what the
predictor contributes from what the Huffman dictionary contributes.

\subsection{The standard two-stage structure}

Almost every practical lossless compressor has the same shape: a
\emph{decorrelation} stage followed by an
\emph{entropy coding} stage.


The \textbf{decorrelation} stage exploits structure. If consecutive samples are
similar, predict each one from its neighbours and transmit only the
\emph{residual}, the difference between prediction and truth. When the
prediction is good, residuals cluster tightly around zero even though the
original samples spanned the full dynamic range. Nothing is lost: given the
residual and the same prediction rule, the decoder reconstructs the original
exactly.

The \textbf{entropy coding} stage exploits the resulting skew in the
distribution. Values that occur often get short codewords; values that occur
rarely get long ones. Huffman coding \cite{huffman1952}\label{lit:huffman:bg} is the classic
construction and produces the optimal prefix-free code for a known symbol
distribution.

Crucially, the two stages are independent. A predictor that fails --- that
produces residuals no more skewed than its input --- does not break the coder;
the coder simply carries on compressing whatever it is handed, and the
compression ratio can still look respectable. Detecting this requires
deliberately measuring what the coder achieves \emph{without} the predictor.
This thesis reports that no-prediction baseline alongside every predictive
result.

\subsection{Open-loop versus closed-loop prediction}
\label{sec:closedloop}

One subtlety in predictive coding matters enough to state now, because the
hardware of Chapter~\ref{ch:hardware} is built around it
(Section~\ref{sec:robust}).

The encoder predicts from past samples. So does the decoder --- but the decoder
does not have the past \emph{original} samples, only the past
\emph{reconstructed} ones. In \emph{open-loop} prediction the encoder uses the
originals; in \emph{closed-loop} prediction (classically, differential
pulse-code modulation, DPCM) the encoder reconstructs each sample exactly as
the decoder will, and predicts from that.

When compression is perfectly lossless the two are identical, because
reconstruction equals the original. The moment anything --- a saturation, a
rounding, a clamp --- makes reconstruction differ from the original by even one
bit, an open-loop encoder and its decoder begin predicting from different
histories and \emph{diverge}. With an IIR predictor, as in DRHE, that
divergence is then carried forward through every subsequent ramp. Closing the
loop confines the error to the affected sample.


\section{The sensor and the dataset}
\label{sec:sensor}

\subsection{AWR2243 and the MMWCAS-RF-EVM}

The radar front end assumed throughout is the Texas Instruments AWR2243, a
76--\SI{81}{\giga\hertz} FMCW transceiver with three transmit and four receive
channels \cite{ti_awr2243}. Importantly for this work, the AWR2243 performs no
signal processing of its own: it transmits, receives, mixes and digitises, then
streams raw ADC samples out over a MIPI CSI-2 interface. It therefore exposes
exactly the payload that creates the transport bottleneck.

Four AWR2243 devices are cascaded on the MMWCAS-RF-EVM evaluation module,
giving \num{12} transmit and \num{16} receive channels on one synchronised
board --- the \num{192}-element virtual array of Section~\ref{sec:mimo}.

\subsection{ColoRadar}

Rather than synthetic stimulus, this thesis uses real captures from the public
ColoRadar dataset \cite{kramer2022}\label{lit:kramer:bg}, which includes cascade recordings from
exactly this module across indoor, outdoor and subterranean environments.
Table~\ref{tab:capture} lists the relevant parameters. Fifty consecutive frames
were used throughout; that is \num{39321600} complex ADC samples, or
\SI{150}{\mebi\byte} of raw data, which is enough to make a claim of
losslessness meaningful.

\begin{table}[t]
\centering
\caption{ColoRadar cascade capture parameters used throughout this thesis.}
\label{tab:capture}
\begin{tabular}{@{}llr@{}}
\toprule
Parameter & Symbol & Value \\
\midrule
ADC samples per chirp      & $N_s$ & 256 \\
Chirp loops per frame      & $N_c$ & 16 \\
Receive channels           & $N_r$ & 16 \\
Transmit channels (TDM)    & $N_t$ & 12 \\
Ramps per frame            & $N_c N_t$ & 192 \\
Fast-time sample rate      & $f_s$ & \SI{8}{\mega\hertz} \\
Chirp bandwidth            & $B$ & \SI{1.798}{\giga\hertz} \\
Ramp repetition interval   & $T_r$ & \SI{78.5}{\micro\second} \\
Centre frequency           & $f_c$ & \SI{77}{\giga\hertz} \\
Range resolution           & $\Delta R$ & \SI{8.34}{\centi\metre} \\
Frame size (raw)           & --- & \SI{3145728}{\byte} \\
\bottomrule
\end{tabular}
\end{table}

Each complex sample is stored as two contiguous \texttt{int16} values. The
index of the sample at (fast-time $s$, chirp loop $q$, receiver $r$,
transmitter $t$) in the file is
\begin{equation}
i = 2\bigl(s + N_s(q + N_c(r + t N_r))\bigr),
\end{equation}
so the layout from slowest to fastest varying is $[t, r, q, s]$. Reshaping this
into the (fast time, ramp, receiver) cube the processing chain expects is what
produces the transmitter-interleaved ramp axis of Equation~\eqref{eq:rampindex}.

\section{Quantisation: the FX16 step}
\label{sec:numbers}
\label{sec:fx16}

The range-FFT output is naturally a complex floating-point quantity. Before it
can be compressed by a fixed-point pipeline it has to be quantised to 16-bit
integers --- the \textbf{FX16} step. This is a scaling followed by rounding:
\begin{equation}
\label{eq:scale}
\sigma = \frac{32767}{A_{\max}\,2^{\mathrm{shift}} \sum_k w_k \,/\, (2 N_s)},
\qquad
x_{\mathrm{FX16}} = \mathrm{round}(\sigma \cdot x),
\end{equation}
where $A_{\max}$ is the full-scale ADC amplitude and $w$ is the window applied
along fast time.

Two things about Equation~\eqref{eq:scale} matter later. First, the rounding
introduces an error bounded by half a least-significant bit, which acts as a
wideband noise floor \cite{widrow2008}; this is the dominant source of
detection degradation in the whole chain, and Chapter~\ref{ch:results}
separates it cleanly from the cost of compression. Second, and more
consequentially for this thesis, the scale factor is \emph{fixed} --- it is
calibrated from a theoretical full-scale ADC sinusoid, not from the data.
If the actual capture sits well below full scale, the top bits of every
resulting \texttt{int16} are only copies of the sign bit. On ColoRadar the largest FX16
code is \num{437}, so fewer than ten of the sixteen bits ever carry
information; Section~\ref{sec:ratiomeaning} shows how much of the compression
ratio this accounts for.


% -------------------------------------------------------------------
\section{FPGAs and high-level synthesis}
\label{sec:fpga}

A field-programmable gate array (FPGA) is a chip of configurable logic blocks
and wiring that \emph{becomes} the circuit described, rather than executing
instructions. It can process four receive channels in the same clock cycle,
with deterministic timing, and can be changed after manufacture, which is why
it suits a proof of concept; a production design would normally move into an
ASIC or a fixed-function block of a larger SoC \cite{kuon2008,nxp_an5375}.
Its resources are counted in four units: \textbf{LUTs} (look-up tables, which
implement logic), \textbf{flip-flops} (one-bit registers), \textbf{block RAM}
(on-chip memory, counted here in \SI{36}{\kilo bit} blocks, each with two
access ports) and \textbf{DSP slices} (hardened multiply--accumulate units).
The device used throughout, an AMD Kintex UltraScale+
\texttt{xcku5p-ffvb676-2-e}, provides \num{216960} LUTs, \num{433920}
flip-flops, 480 block RAMs and \num{1824} DSP slices. Fixed-point arithmetic
maps onto these cheaply; floating point needs a multi-cycle operator for every
operation \cite{parhami2010,ieee754}.

\label{sec:pipelining}
A \textbf{pipeline} splits a computation into stages separated by registers so
that several inputs are in flight at once; its \emph{depth} is the number of
cycles one input takes to pass through. The \textbf{initiation interval} (II)
is the number of cycles between accepting one input and the next:
$\mathrm{II} = 1$ is one input per clock, and $\mathrm{II} = 2$ halves the
throughput. $\mathbf{F_{\max}}$ is the fastest clock at which every path still
settles within one period. A frame of $K$ input beats takes about
$K \cdot \mathrm{II} / f$ seconds, plus the pipeline depth once per loop
invocation.

\textbf{High-level synthesis} (HLS) generates the circuit from annotated C++:
loops become pipelined state machines, arrays become memories, and pragmas set
the II and the interfaces. Because the source is C++, the same code also runs
as software against the MATLAB reference before any hardware exists, which
makes the verification of Section~\ref{sec:ladder} possible. The cores here use
AXI4-Stream interfaces: a valid/ready handshake carrying one fixed-width word
per beat.

% -------------------------------------------------------------------
\section{Related work}
\label{sec:sota}
\label{sec:kiem}

Commercial radar processors use lightweight deterministic compressors with
bounded latency: Texas Instruments' exponential-Golomb encoding \cite{ti_ege}
and NXP's block floating-point formats such as CP4D \cite{nxp_an5375}, both
``near lossless'' at ratios of about 2 to 4. Block adaptive quantisation
(BAQ), from synthetic aperture radar \cite{kwok1989}, adapts the quantiser
step to each block's statistics. Range-dependent variable-length encoding
(RDVLE) \cite{li2014} spends fewer bits on far range bins, where received power
is lower. All three are lossy. Doppler Redundancy Hybrid Encoding (DRHE)
\cite{li2014,zivkovic2015} predicts each range-FFT sample along slow time with
two IIR filters in polar form and entropy-codes the error; linear predictive
coding (LPC) is the general form of the same idea, applied by Meucci and
Mancuso \cite{meucci2022} to raw ADC samples. Autoencoder and spectral
approaches reach higher ratios but are lossy and expensive at the sensor
\cite{cheng2021,park2026}.

The direct predecessor of this work is Lukas Kiem's master's thesis
\cite{kiem2025}, which compared industrial schemes with DRHE on a
single-transmitter Infineon radar, selected DRHE (3.25 in simulation) and
implemented it on a Zynq UltraScale+ device (3.17 in hardware).
Chapter~\ref{ch:kiem} summarises it.
